\section{Erd\H{o}s Problem 369}

\paragraph{Problem statement.}
Let $\varepsilon>0$ and $k\geq 2$. Is it true that, for all sufficiently large $n$, there is a sequence of $k$ consecutive integers in $\{1,\ldots,n\}$ all of which are $n^{\varepsilon}$-smooth?

\paragraph{Tools in this round.}
Web browsing: YES. Computation: YES.

\subsection*{1) ROUND-2 OBJECTIVE}
Path (C): obstruction/correction. Round 1 already gave a complete proof of the literal statement, so the only mathematically meaningful Round-2 task is to identify the exact defect in the formulation and determine which corrected form actually captures a nontrivial asymptotic problem.

The plan is to separate two possible repairs:
\begin{itemize}
\item replace ``$n^{\varepsilon}$-smooth'' by ``$m^{\varepsilon}$-smooth'' for each term $m$ of the block;
\item force the block itself to lie near the top of the interval, for example inside $[n/2,n]$.
\end{itemize}
I will prove that the first repair is true (indeed much weaker than it first appears), whereas the second repair is the one that preserves the intended asymptotic content and remains open.

\subsection*{2) Round-1 FOUNDATION USED}
I use the following Round-1 results.
\begin{itemize}
\item Round 1 proved the literal statement exactly as written by taking the fixed block $1,2,\ldots,k$ once $n\geq \max\{k,k^{1/\varepsilon}\}$.
\item Round 1 identified the core ambiguity: the phrase ``$n^{\varepsilon}$-smooth'' ties the smoothness threshold to the ambient parameter $n$, not to the size or location of the block itself.
\item Round 1 also noted two natural repairs already discussed in the literature: termwise $m^{\varepsilon}$-smoothness, and a localization condition such as placing the block in $[n/2,n]$.
\end{itemize}

\subsection*{3) NEW INSIGHT / TOOL (ROUND-2)}
The new point in this round is the following adversarial observation.

Merely replacing $n^{\varepsilon}$-smooth by $m^{\varepsilon}$-smooth does \emph{not} restore genuine asymptotic dependence on $n$. Once a single admissible block exists anywhere, the same block lies in $\{1,\ldots,n\}$ for every larger $n$. So this repaired statement is equivalent simply to the existence of one admissible block.

This means that the true logical defect in the original statement is not only the global smoothness threshold; it is also the absence of any location condition. Balog--Wooley's theorem then settles the ``termwise-smoothness'' repair immediately, while the localization repair remains open.

\subsection*{4) ATTACK PLAN (ROUND-2)}
The precise gaps left after Round 1 are:
\begin{itemize}
\item the corrected variants were only identified informally, not resolved;
\item it was not yet shown which corrected variant preserves the asymptotic flavor of the original question.
\end{itemize}

I will now prove three claims.
\begin{enumerate}
\item The termwise correction
\[
(\forall X \text{ large})\; \exists \text{ a block of $k$ consecutive integers in } \{1,\ldots,X\}
\]
whose members $m$ all satisfy $P^{+}(m)\leq m^{\varepsilon}$ is equivalent to the existence of at least one such block.
\item This corrected statement is true, by Balog--Wooley.
\item Therefore the only repair that keeps the parameter $n$ genuinely relevant is a localization requirement such as $P\subseteq [n/2,n]$; that version remains open.
\end{enumerate}

\subsection*{5) WORK (ROUND-2)}
Write $P^{+}(m)$ for the largest prime factor of $m$.

\paragraph{Step 1: the smoothness-only repair collapses to a one-time existence statement.}
Consider the statement
\[
\mathcal C_{\mathrm{smooth}}(\varepsilon,k):\qquad
\text{for all sufficiently large } X, \text{ there exists a block of } k
\text{ consecutive integers in } \{1,\ldots,X\}
\]
all of whose members $m$ satisfy $P^{+}(m)\leq m^{\varepsilon}$.

I claim that $\mathcal C_{\mathrm{smooth}}(\varepsilon,k)$ is equivalent to the simpler assertion:
\[
\mathcal E(\varepsilon,k):\qquad
\text{there exists at least one block of } k \text{ consecutive positive integers }
\]
all of whose members $m$ satisfy $P^{+}(m)\leq m^{\varepsilon}$.

Indeed, if $\mathcal C_{\mathrm{smooth}}(\varepsilon,k)$ holds then certainly $\mathcal E(\varepsilon,k)$ holds. Conversely, suppose one block
\[
M+1,M+2,\ldots,M+k
\]
satisfies $P^{+}(M+j)\leq (M+j)^{\varepsilon}$ for every $1\leq j\leq k$. Then the same block is contained in $\{1,\ldots,X\}$ for every $X\geq M+k$. Hence $\mathcal C_{\mathrm{smooth}}(\varepsilon,k)$ holds.

So the ``for all sufficiently large $X$'' quantifier is not the real issue here: without a location constraint, one admissible block is enough forever.

\paragraph{Step 2: prove $\mathcal E(\varepsilon,k)$.}
There are two cases.

\smallskip
\noindent\emph{Case 1: $\varepsilon\geq 1$.}
For every positive integer $m$ we have $P^{+}(m)\leq m\leq m^{\varepsilon}$. Therefore every block of consecutive integers works; in particular $1,2,\ldots,k$ is admissible.

\smallskip
\noindent\emph{Case 2: $0<\varepsilon<1$.}
Put $u=1/\varepsilon$, so $u>1$. Balog--Wooley prove that if
\[
t(N)=\left\lfloor \frac{\log_{4} N}{\log(3u)}\right\rfloor,
\]
then for infinitely many positive integers $N$, none of the prime factors of any of the integers
\[
N+1,N+2,\ldots,N+t(N)
\]
exceeds $N^{1/u}=N^{\varepsilon}$.

Since $t(N)\to\infty$, there exists such an $N$ with $t(N)\geq k$. Fix one. Then for every $1\leq j\leq k$,
\[
P^{+}(N+j)\leq N^{\varepsilon}< (N+j)^{\varepsilon},
\]
because $\varepsilon>0$ and $N+j>N$. Hence the block
\[
N+1,N+2,\ldots,N+k
\]
satisfies $P^{+}(m)\leq m^{\varepsilon}$ termwise. Thus $\mathcal E(\varepsilon,k)$ holds.

By Step 1, $\mathcal C_{\mathrm{smooth}}(\varepsilon,k)$ holds as well.

\paragraph{Step 3: identify the genuinely asymptotic correction.}
Now consider instead the localized variant
\[
\mathcal C_{\mathrm{loc}}(\varepsilon,k):\qquad
\text{for all sufficiently large } X,\text{ there exists a block of } k \text{ consecutive integers }
\]
contained in $[X/2,X]$ all of which are $X^{\varepsilon}$-smooth.

The monotonicity trick from Step 1 no longer applies: an early block does \emph{not} stay inside $[X/2,X]$ as $X\to\infty$. Thus the parameter $X$ again becomes essential. This is therefore the correction that preserves the intended asymptotic content. According to the current literature summary checked in this round, this localized version is open for ``all sufficiently large $X$''.

\subsection*{6) ADVERSARIAL VERIFICATION}
I check the vulnerable points explicitly.
\begin{itemize}
\item \emph{Did the termwise repair really collapse to one-time existence?} Yes. The repaired statement asks only for a block somewhere in the initial segment $\{1,\ldots,X\}$. Once one admissible block exists, every larger initial segment still contains it.
\item \emph{Did I use monotonicity correctly when $0<\varepsilon<1$?} Yes. The function $x\mapsto x^{\varepsilon}$ is increasing for every $\varepsilon>0$, so $N^{\varepsilon}<(N+j)^{\varepsilon}$.
\item \emph{Does Balog--Wooley give length exactly $k$?} Not exactly, but it gives length $t(N)$ with $t(N)\to\infty$, so for fixed $k$ there are admissible strings of length at least $k$, and the first $k$ terms give the desired block.
\item \emph{Is the localization correction really immune to the Round-1 trivialization?} Yes. A block near the origin does not remain in $[X/2,X]$ once $X$ grows, so the quantifier ``for all sufficiently large $X$'' regains genuine content.
\end{itemize}
No contradiction with Round 1 appears; the new work sharpens the logical diagnosis rather than undoing any previous lemma.

\subsection*{7) FINAL}
\textbf{FULL SOLUTION -- FULL PROOF.}

The statement as written is fatally underconstrained and is trivially true, exactly as Round 1 showed. In Round 2 I proved more:
\begin{itemize}
\item the smoothness-only repair (replace $n^{\varepsilon}$-smooth by $m^{\varepsilon}$-smooth) is also true, because it is equivalent to the existence of one admissible block and Balog--Wooley provide such blocks;
\item the asymptotically meaningful repair is the localized one (for example, requiring the block to lie in $[n/2,n]$), and that is the version that remains open.
\end{itemize}
Thus the original problem cannot be the intended hard problem in its literal form.

\subsection*{8) COMPLETION ESTIMATE (MANDATORY)}
\textbf{COMPLETION: 100\%}

\subsection*{9) REFERENCES}
\begin{itemize}
\item A. Balog and T. D. Wooley, \emph{On strings of consecutive integers with no large prime factors}, J. Austral. Math. Soc. Ser. A 64 (1998), 266--276.
\item T. F. Bloom, \emph{Erd\H{o}s Problem \#369}, accessed 2026-03-07.
\end{itemize}

\subsection*{10) OUTPUT FORMAT}
Erd\H{o}s problem number: $x=369$. This section is written as a drop-in \LaTeX{} fragment and is included in the combined file \texttt{369\_371\_373\_round2\_solutions.tex}.

\bigskip
\hrule
\bigskip

